\chapter{Kết luận và hướng nghiên cứu tương lai}
\label{Chapter5}

Trong báo cáo này, chúng tôi đã đề xuất hai mô-đun toán tử truy vấn liên tục và rời rạc, cùng phiên bản ngẫu nhiên tương ứng với từng toán tử phục vụ việc định lượng độ bất định. Ngoài ra, chúng tôi đề xuất một phương pháp thích ứng khi kiểm thử sử dụng cơ chế phân luồng và hàm mất mát tương phản giữa nhánh biểu diễn trực tuyến và ngoại tuyến. Kết quả thực nghiệm cho thấy pha ngoại tuyến và pha trực tuyến của phương pháp đều tốt hơn so với các phương pháp đối sánh được chúng tôi lựa chọn. Trong pha trực tuyến, việc sử dụng cơ chế làm trơn đóng góp nhiều nhất vào mức tăng hiệu quả. Kết quả định lượng bất định cũng cho thấy mạng có độ bất định cao hơn khi tính toán trên tập kiểm tra và các độ đo vẫn cho thấy mạng tổng quát hoá tốt trên dữ liệu chưa được học khi huấn luyện. Cuối cùng, chúng tôi xin đề xuất hai hướng nghiên cứu tương lai như sau:
1. Giảm \textit{tỉ lệ báo động giả (false positive rate)} để đẩy lùi thực trạng \textit{mệt mỏi vì báo động (alarm fatigue)} trong một số lĩnh vực như giám sát hệ thống phần mềm và đám mây, hay trong các đơn vị hồi sức tích cực, và 
2. Nghiên cứu cách đánh giá chất lượng ước lượng độ bất định để có cơ sở so sánh công bằng các phương pháp ước lượng độ bất định khác nhau trong bài toán phát hiện bất thường chuỗi thời gian đa biến.